% Source: https://artofproblemsolving.com/community/c3h2686197_function_transformation
% Stable identity: aops-topic-2686197-q01
\textbf{(b)} A point $(u,v)$ is on the graph of $g$ exactly when $v=2f(u-6)-3$, that is, when $\dfrac{v+3}{2}=f(u-6)$, i.e. when $\left(u-6,\dfrac{v+3}{2}\right)$ is on the graph of $f$. Writing $x=u-6$ and $y=\dfrac{v+3}{2}$ gives $u=x+6$ and $v=2y-3$. So the graph of $g$ is obtained from the graph of $f$ by
shifting right $6$, then stretching vertically by the factor $2$ (away from the $x$-axis), then shifting down $3$. The stretch must come before the downward shift.

\textbf{(a)} By (b), $(4,5)$ goes to $(4+6,\,2\cdot5-3)=\boxed{(10,7)}$. Check: $g(10)=2f(4)-3=7$.

\textbf{(c)} \textbf{No.} Shifting down $3$ first and then stretching sends $y$ to $2(y-3)=2y-6$, which gives the graph of $2f(x-6)-6$, not $g$.
